\documentclass[12pt]{article}
\usepackage[utf8]{inputenc}
\usepackage[T1]{fontenc}
\usepackage{hyperref}
\title{Sicula Language Resources for K-12 Curriculum Development}
\author{Compiled by Claude Sonnet 5.5 (claude-sonnet-5-5) for Andr\'as Kornai}
\date{}
\begin{document}
\maketitle

\section{Introduction}
Sicula (Siculian, Sicel; ISO 639-3 scx) is an extinct language of central and eastern Sicily, attested in fewer than thirty inscriptions from the late 6th to the 4th century BCE and about twenty-five glosses in ancient authors. Most modern scholars favour an Italic affiliation, but the evidence does not allow a firm conclusion. Markus Hartmann judged that only about six inscriptions are at least partly intelligible. There are no native speakers, so a curriculum ``in the language'' is not possible. Only study of the language and its culture is possible.

\section{Primary Educational Resources}
I found no textbooks or school materials. The surviving texts are too few for language teaching, though they could support classroom units on ancient Sicily, epigraphy and writing systems.

\section{Linguistic and Reference Materials}
\begin{itemize}
\item Wikipedia, Siculian (summary of the corpus, script and classification debate): \url{https://en.wikipedia.org/wiki/Siculian}
\item Wikipedia, Sicels (the people): \url{https://en.wikipedia.org/wiki/Sicels}
\item I.Sicily, digital corpus of Sicilian inscriptions (includes Sikel texts), project page: \url{https://dh.web.ox.ac.uk/project/isicily-digital-corpus-sicilian-inscriptions}
\item I.Sicily entry ISic003362, a stele in the Sikel language from Licodia Eubea: \url{https://isicily.classics.ox.ac.uk/inscription/ISic003362}
\end{itemize}

\section{Religious and Community Resources}
None found. There is no living community that uses the language.

\section{Digital and Audio Resources}
The I.Sicily online corpus is the main digital resource (stone inscriptions only). No audio resources were found.

\section{Curriculum Development Context}
Only secondary summaries and the I.Sicily corpus were consulted. Any unit should draw on specialist scholarship, such as Hartmann's work cited by Wikipedia, and should say plainly that the language is poorly understood.
\end{document}